% !TEX root = ../main.tex
% Figure: problem-first method map for PCSEL simulation

\begin{figure}[htbp]
  \centering
  \resizebox{\textwidth}{!}{%
  \begin{tikzpicture}[x=1cm,y=1cm,>=Latex,font=\small]
    \tikzset{
      question/.style={draw=black!60,rounded corners=2pt,fill=gray!6,align=center,minimum width=3.2cm,minimum height=0.78cm,font=\footnotesize},
      method/.style={draw=blue!55!black,rounded corners=2pt,fill=blue!7,align=center,minimum width=2.75cm,minimum height=0.82cm,font=\bfseries\footnotesize},
      output/.style={draw=green!45!black,rounded corners=2pt,fill=green!7,align=center,text width=3.15cm,minimum height=1.05cm,font=\footnotesize},
      caution/.style={draw=red!60!black,rounded corners=2pt,fill=red!6,align=left,text width=4.45cm,font=\footnotesize},
      flow/.style={->,thick},
      softflow/.style={->,thick,dashed,gray!75}
    }

    \node[font=\bfseries] at (7.5,5.25) {先按问题类型选方法，而不是按软件名字倒推结论};

    \node[question] (q1) at (1.8,4.1) {候选频段在哪里?};
    \node[method] (m1) at (1.8,2.95) {PWEM / 能带};
    \node[output] (o1) at (1.8,1.55) {频率窗口\\简并关系\\高对称点标签};

    \node[question] (q2) at (5.6,4.1) {主导耦合是谁?};
    \node[method] (m2) at (5.6,2.95) {CWT / CMT};
    \node[output] (o2) at (5.6,1.55) {散射通道\\参数灵敏度\\模式分裂方向};

    \node[question] (q3) at (9.4,4.1) {开放单胞如何辐射?};
    \node[method] (m3) at (9.4,2.95) {RCWA};
    \node[output] (o3) at (9.4,1.55) {上下出光\\辐射损耗\\偏振趋势};

    \node[question] (q4) at (13.2,4.1) {有限器件是否仍成立?};
    \node[method] (m4) at (13.2,2.95) {FDTD / FEM};
    \node[output] (o4) at (13.2,1.55) {有限阵列近远场\\边缘损耗\\真实层栈修正};

    \foreach \a/\b in {q1/m1,m1/o1,q2/m2,m2/o2,q3/m3,m3/o3,q4/m4,m4/o4} {
      \draw[flow] (\a) -- (\b);
    }
    \foreach \a/\b in {o1/q2,o2/q3,o3/q4} {
      \draw[softflow] (\a.east) -- ++(0.35,0) |- (\b.west);
    }

    \node[caution,anchor=north west] at (0.25,0.25) {
      \textbf{层级边界：}\\
      PWEM/CWT 不能单独给阈值电流；\\
      RCWA 不能单独给有限阵列边缘排序；\\
      冷腔 FDTD/FEM 不能单独给热-电-光工作窗口。
    };
    \node[draw=black!45,rounded corners=2pt,fill=yellow!10,align=center,text width=8.6cm,font=\footnotesize,anchor=north west] at (5.35,0.25) {
      最可靠的结论通常来自“同一物理解释链”在多个层级上保持一致：\\
      候选模身份稳定、对称性标签稳定、参数趋势稳定、损耗来源可追踪。
    };
  \end{tikzpicture}%
  }
  \caption{PCSEL 数值方法的分工地图。图中强调从问题类型进入：能带/PWEM 给候选目录，CWT/CMT 给机理压缩，RCWA 给开放周期单胞辐射，FDTD/FEM 给有限器件与真实层栈验证。它是教学示意图，不代表跨软件 benchmark 分数。}
  \label{fig:ch15_method_comparison_radar}
\end{figure}
